\documentclass[a4paper,10pt,twocolumn]{ctexart}

\usepackage[a4paper,margin=1.7cm,columnsep=0.7cm]{geometry}
\usepackage{graphicx}
\usepackage{booktabs}
\usepackage{array}
\usepackage{multirow}
\usepackage{amsmath}
\usepackage{url}
\usepackage{hyperref}
\usepackage{float}
\usepackage{xcolor}
\usepackage{caption}
\usepackage{enumitem}
\usepackage{fancyhdr}

\hypersetup{
    colorlinks=true,
    linkcolor=black,
    citecolor=black,
    urlcolor=blue
}

\newcommand{\reporttitle}{利用深度学习框架对比微调和从头训练}
\newcommand{\coursename}{人工神经网络}
\newcommand{\assignmentname}{报告二}
\newcommand{\studentname}{待填写}
\newcommand{\studentid}{待填写}
\newcommand{\teachername}{待填写}
\newcommand{\repourl}{https://github.com/yourname/ann-hw2-finetune-vs-scratch}

\newcommand{\BestAccModel}{ResNeXt-FT}
\newcommand{\BestAccValue}{79.69\%}
\newcommand{\BestFOneModel}{ResNeXt-FT}
\newcommand{\BestFOneValue}{79.53\%}
\newcommand{\FastestModel}{DenseNet-FT}
\newcommand{\FastestTime}{5.6}
\newcommand{\SlowestModel}{ResNeXt-FT}
\newcommand{\SlowestTime}{6.8}


\ctexset{
    section={
        format=\Large\bfseries,
        number=\arabic{section},
        beforeskip=0.8ex,
        afterskip=0.6ex
    },
    subsection={
        format=\normalsize\bfseries,
        beforeskip=0.5ex,
        afterskip=0.4ex
    }
}

\pagestyle{fancy}
\fancyhf{}
\fancyhead[C]{\coursename\quad \assignmentname}
\fancyfoot[C]{\thepage}
\setlength{\headheight}{14pt}
\setlength{\parindent}{2em}
\setlength{\parskip}{0.2em}
\captionsetup{font=small}
\newcommand{\resultfigure}[4]{
    \begin{figure}[t]
        \centering
        \IfFileExists{#1}{
            \includegraphics[width=#2]{#1}
        }{
            \fbox{\parbox[c][4.0cm][c]{0.9\columnwidth}{\centering 图像将在实验完成后由脚本自动生成}}
        }
        \caption{#3}
        \label{#4}
    \end{figure}
}

\begin{document}

\twocolumn[
\begin{center}
    {\LARGE \bfseries \reporttitle \par}
    \vspace{0.8em}
    {\normalsize 课程：\coursename \quad 作业：\assignmentname \par}
    \vspace{0.4em}
    {\normalsize 姓名：\studentname \quad 学号：\studentid \quad 任课教师：\teachername \par}
    \vspace{0.8em}
\end{center}

\begin{quote}
\small
\textbf{摘要：}
本文围绕卷积神经网络在小样本图像分类任务中的训练策略展开比较，重点分析从头训练与利用公开预训练权重进行微调两种方案的差异。实验选择 ResNeXt50\_32x4d 和 DenseNet121 两种代表性 CNN 模型，在 Beans 三分类数据集上构建四组对照实验，并从 Accuracy、Macro-F1、训练时间和收敛曲线等角度进行分析。结果表明，在本实验设置下，微调方案整体上具有更快的收敛速度和更稳定的测试表现；与此同时，不同模型结构在小规模数据集上的训练稳定性也存在差异。本文最后结合实际训练过程，总结了小样本场景下模型选择和训练策略上的一些体会。
\end{quote}
\vspace{0.8em}
]

\section{问题简介}

图像分类是计算机视觉中的基础任务，而在真实应用中，研究者经常会遇到一个很实际的问题：目标数据集规模较小，但模型结构已经较成熟，此时究竟应该从头训练，还是直接利用公开预训练权重进行微调。这个问题在课程实验中也很有代表性，因为它不仅关系到最终分类精度，还会影响训练时间、参数收敛速度以及调参难度。

本次实验围绕这一问题展开。本文选择两种典型的卷积神经网络结构，即 ResNeXt50\_32x4d 和 DenseNet121，在 Beans 三分类数据集上分别进行从头训练和微调，共形成四组对照实验。这样的设置有两个考虑：一方面，ResNeXt 与 DenseNet 在特征提取方式上存在明显差异，前者强调 grouped transformation 与 cardinality，后者强调特征复用和密集连接；另一方面，小规模农业病害数据集也比较适合观察迁移学习在有限样本条件下的效果。

本实验希望回答三个问题。第一，在小样本图像分类任务中，微调是否一定优于从头训练。第二，不同 CNN 结构在同一数据集上的表现是否一致。第三，在 CPU 条件下完成训练时，哪一种训练策略更稳定、更容易复现。从课程作业的角度看，这三个问题既有理论比较意义，也与实际计算成本密切相关。

\section{模型与算法介绍}

\subsection{ResNeXt 与 DenseNet}

ResNeXt 是在残差网络基础上发展起来的一类模型，其核心思想是将多条具有相同拓扑结构的变换分支进行聚合，并通过 cardinality 提升表示能力。与单纯增加网络深度或宽度相比，这种做法在保持结构规则性的同时，能够更高效地提升特征表达能力。本文选用的 ResNeXt50\_32x4d 是一个较常见的标准版本，适合与其他经典 CNN 进行对比。

DenseNet 的特点是密集连接，即每一层都会接收前面所有层的特征图作为输入，并将自己的输出继续传递给后续层。这样的设计有利于特征复用，也能够在一定程度上缓解梯度传播过程中的信息损失。本文选用 DenseNet121 作为第二个对照模型，希望观察其在小样本图像分类中的稳定性表现。

\subsection{从头训练与微调}

从头训练是指模型参数全部随机初始化，然后仅依靠当前任务数据进行优化。这种方式的优点是训练过程完全针对目标任务展开，不依赖外部数据；但缺点也比较明显，当数据量较小时，模型往往更依赖超参数设置，训练初期容易震荡，收敛速度通常较慢。

微调则是将模型初始化为在 ImageNet 上训练得到的公开预训练权重，再针对当前任务继续训练。考虑到预训练模型已经学到了较通用的边缘、纹理和中层语义特征，微调通常会在小样本任务中表现出更快的收敛速度和更高的起始性能。为了让对比更公平，本文对两种模型都采用统一的两阶段微调方案：先冻结主干网络，仅训练最后的分类层；随后解冻全部参数，以较小学习率进行全模型优化。

\subsection{训练设置}

四组实验均使用交叉熵损失函数和 AdamW 优化器，输入图像统一缩放到 $224 \times 224$。从头训练的初始学习率设置为 $3 \times 10^{-4}$，训练 30 个 epoch；微调实验先以 $1 \times 10^{-3}$ 训练分类头 5 个 epoch，再以 $1 \times 10^{-4}$ 解冻全模型训练 15 个 epoch。评价指标使用测试集 Accuracy 和 Macro-F1，并记录训练总时长、最优轮次以及训练过程中的曲线变化。

\section{数据集和算例参数介绍}

\subsection{数据集说明}

本文采用 Beans 数据集进行实验。该数据集来自豆类叶片病害识别任务，包含 \texttt{angular\_leaf\_spot}、\texttt{bean\_rust} 和 \texttt{healthy} 三个类别，具有比较明确的农业图像分类背景。数据集官方已经给出了训练集、验证集和测试集划分，因此本文直接沿用原始划分，以避免额外的数据切分对对比结果造成影响。

表~\ref{tab:dataset} 给出了数据划分情况。可以看出，Beans 数据集整体规模不大，这也是本文重点比较从头训练和微调的原因之一：在样本较少时，预训练特征往往更容易体现优势。

\begin{table}[t]
    \caption{Beans 数据集划分统计}
    \label{tab:dataset}
    \centering
    \resizebox{\columnwidth}{!}{\begin{tabular}{p{0.17\columnwidth}p{0.12\columnwidth}p{0.52\columnwidth}}
\toprule
数据划分 & 样本数 & 类别说明 \\
\midrule
Train & 1034 & 角斑病, 锈病, 健康叶片 \\
Validation & 133 & 三分类 \\
Test & 128 & 三分类 \\
\bottomrule
\end{tabular}
}
\end{table}

\subsection{预处理与增强}

训练阶段采用 RandomResizedCrop、随机水平翻转和轻量级 ColorJitter，验证和测试阶段使用 Resize 与 CenterCrop。图像归一化采用 ImageNet 常用的均值和标准差，这样做的原因有两点：一是方便与预训练模型保持一致，二是保证 scratch 和 fine-tune 两种方案在输入分布处理上尽可能统一。

\subsection{硬件环境与记录方式}

实验在 CPU 环境下完成，框架为 Python、PyTorch 和 torchvision。由于课程作业强调训练过程和实验比较，因此除了最终准确率外，本文还记录了每组实验的训练总时长、最优验证轮次以及训练过程中的 loss/accuracy 曲线。训练时间采用脚本内部的 wall-clock 统计方式，覆盖完整训练流程，便于在相同硬件条件下直接比较不同策略的成本。

\section{结果罗列与汇总}

\subsection{整体结果比较}

四组实验的主要结果见表~\ref{tab:main-results}。从整体上看，在本实验设置下，\BestAccModel{} 取得了最高的测试集准确率，为 \BestAccValue{}；而在 Macro-F1 指标上，\BestFOneModel{} 的表现最好，达到 \BestFOneValue{}。这一结果说明，在样本规模较小且类别较少的图像分类任务中，利用预训练特征通常更容易获得较稳定的分类边界。

\begin{table}[t]
    \caption{四组实验的测试集结果对比}
    \label{tab:main-results}
    \centering
    \resizebox{\columnwidth}{!}{\begin{tabular}{lcccc}
\toprule
实验组别 & Acc(\%) $\uparrow$ & Macro-F1 $\uparrow$ & 最优轮次 & 参数量(M) \\
\midrule
ResNeXt-S & -- & -- & -- & -- \\
ResNeXt-FT & 79.69 & 79.53 & 1 & 22.99 \\
DenseNet-S & -- & -- & -- & -- \\
DenseNet-FT & 75.78 & 75.91 & 1 & 6.96 \\
\bottomrule
\end{tabular}
}
\end{table}

从训练方式看，微调方案在训练前期通常能更快进入有效收敛区间，而从头训练的初始波动更明显。这一点在实际训练时体会比较直接：如果学习率稍大，随机初始化模型在前几个 epoch 容易出现验证集指标反复波动，而使用预训练权重的模型往往在较早阶段就能得到较高的验证分数。对于 Beans 这样的数据集，这种差异会比较明显。

\subsection{训练时间比较}

表~\ref{tab:time-results} 给出了不同实验组的训练时间统计。就本次实验而言，训练时间最短的模型是 \FastestModel{}，约为 \FastestTime{} 分钟；训练时间最长的是 \SlowestModel{}，约为 \SlowestTime{} 分钟。这里可以看到，训练成本不仅与是否采用微调有关，也与模型结构本身相关。参数规模较大的网络在 CPU 环境下会明显拉长训练时间，因此在课程作业中同时比较精度和时间是有必要的。

\begin{table}[t]
    \caption{四组实验的训练时间统计}
    \label{tab:time-results}
    \centering
    \resizebox{\columnwidth}{!}{\begin{tabular}{lccc}
\toprule
实验组别 & 训练时长(min) & 权重初始化 & 训练方式 \\
\midrule
ResNeXt-S & -- & -- & -- \\
ResNeXt-FT & 6.8 & ImageNet & 微调 \\
DenseNet-S & -- & -- & -- \\
DenseNet-FT & 5.6 & ImageNet & 微调 \\
\bottomrule
\end{tabular}
}
\end{table}

\subsection{训练曲线与现象分析}

图~\ref{fig:curves-resnext} 和图~\ref{fig:curves-densenet} 展示了两类模型在不同训练方式下的曲线。结合这些曲线可以观察到三个比较典型的现象。第一，微调模型在前几轮的下降速度更快，说明预训练特征确实提供了较好的初始化。第二，从头训练更依赖超参数设置，尤其是在小数据集上，如果增强策略或学习率不合适，验证曲线容易出现起伏。第三，不同模型结构的稳定性并不完全一样，DenseNet 在一些小样本任务上往往表现得更平稳，而 ResNeXt 可能需要更长的时间来发挥结构优势。

\resultfigure{generated/resnext50_scratch_curve.png}{\columnwidth}{ResNeXt50 从头训练的曲线}{fig:curves-resnext-scratch}

\resultfigure{generated/resnext50_finetune_curve.png}{\columnwidth}{ResNeXt50 微调的曲线}{fig:curves-resnext}

\resultfigure{generated/densenet121_scratch_curve.png}{\columnwidth}{DenseNet121 从头训练的曲线}{fig:curves-densenet-scratch}

\resultfigure{generated/densenet121_finetune_curve.png}{\columnwidth}{DenseNet121 微调的曲线}{fig:curves-densenet}

\subsection{混淆矩阵观察}

为了进一步观察模型对三类样本的区分情况，本文还绘制了微调模型在测试集上的混淆矩阵，如图~\ref{fig:cm-resnext} 和图~\ref{fig:cm-densenet} 所示。如果某一类病害样本在纹理和颜色上比较接近，模型就更容易出现相互混淆。结合混淆矩阵可以更具体地看出模型误差主要集中在哪些类别之间，而不是只看整体准确率。

\resultfigure{generated/resnext50_finetune_cm.png}{0.92\columnwidth}{ResNeXt50 微调模型的混淆矩阵}{fig:cm-resnext}

\resultfigure{generated/densenet121_finetune_cm.png}{0.92\columnwidth}{DenseNet121 微调模型的混淆矩阵}{fig:cm-densenet}

\section{结论或计算心得}

本文在 Beans 三分类任务上，对 ResNeXt50\_32x4d 和 DenseNet121 的从头训练与微调方案进行了比较。综合 Accuracy、Macro-F1 与训练时间来看，微调在本实验设置下整体更有优势，尤其是在训练初期表现出更快的收敛速度。这说明在小样本图像分类任务中，ImageNet 预训练权重所提供的通用特征仍然具有较强的迁移价值。

从训练过程中的实际体会来看，从头训练并不是不能取得较好的结果，但它对超参数更敏感，训练前期也更容易出现震荡。如果只看最终结果，可能会觉得两种方法差别不算特别大；但真正跑完整个实验后，会发现微调不仅更容易达到较高指标，而且调参压力也相对更小。尤其是在 CPU 环境下，重复试验的成本比较高，这时稳定性的重要性会更加突出。

同时，本实验也说明模型结构的差异会影响训练表现。DenseNet 依靠特征复用，在小规模数据集上往往更容易表现出稳定性；ResNeXt 的结构表达能力较强，但在训练成本和收敛节奏上也可能更依赖具体设置。对课程实验而言，这一点给我的启发是：选择模型时不能只看“结构是否更先进”，还需要结合数据规模、硬件条件以及调参成本一起考虑。

本次实验仍有一定局限。例如，数据集规模较小，且每组实验只使用单次随机种子结果，没有继续做更充分的重复实验；另外，实验环境为 CPU，训练时间相对较长，也限制了更大规模的消融尝试。后续如果条件允许，可以增加更多数据集或加入更多 CNN 结构，进一步分析不同任务中迁移学习收益的变化情况。


\section{参考文献}
\nocite{*}
\bibliographystyle{unsrt}
\bibliography{refs}

\vspace{0.5em}
\noindent\textbf{代码仓库地址：}\url{\repourl}

\end{document}
